\chapter{Gelombang Mekanik}\label{ch:g12:mechanical-waves}

Sabtu malam, sebuah stadion yang penuh: riak penonton yang berdiri
menyapu tribunnya dalam waktu kurang dari satu menit, namun tidak
seorang pun meninggalkan tempat duduknya. Lingkaran riak sebuah kerikil
menyeberangi kolam sementara gabusnya hanya naik turun. Bab ini menelaah
si pengembara itu sendiri --- \omterm{def:g10:signals-and-waves:wave}{gelombang} --- beserta empat bilangan yang
menjinakkannya: laju, \omterm{def:g12:mechanical-waves:delay}{waktu tunda}, \omterm{def:g10:signals-and-waves:period}{periode}, \omterm{def:g12:mechanical-waves:wavelength}{panjang gelombang} --- yang
sudah cukup untuk membaca sebuah gempa dari selembar kertas.

\section{Gangguannya merambat, mediumnya tinggal}

\begin{definition}[Gelombang mekanik]\label{def:g12:mechanical-waves:wave}
\emph{Gelombang mekanik}\index{gelombang mekanik} adalah gangguan yang
merambat menembus medium bendawi --- tali, air, udara, batuan --- sambil
membawa \omterm{def:g10:energy-conservation:energy}{energi} \emph{tanpa mengangkut materi}: tiap potongan mediumnya
berayun sejenak di tempatnya, menyerahkan geraknya, lalu tenang kembali.
\end{definition}

\begin{example}[Tanpa pengangkutan materi]\label{ex:g12:mechanical-waves:transport}
Pada \omterm{def:g10:signals-and-waves:wave}{gelombang} stadion, ``mediumnya'' adalah kerumunan penontonnya: yang
merambat adalah tindakan berdiri, sementara tiap penonton tetap di
tempat duduknya. Di kolam, lingkarannya melebar tetapi gabus yang
mengapung hanya naik turun di tempatnya.
\end{example}

\begin{definition}[Gelombang transversal dan longitudinal]\label{def:g12:mechanical-waves:transverse}
Sebuah \omterm{def:g10:signals-and-waves:wave}{gelombang} disebut \emph{transversal}\index{gelombang transversal}
apabila mediumnya bergerak \emph{tegak lurus} arah rambatnya: sebuah
gundukan melesat menyusuri tali yang disentakkan sementara tiap titiknya
hanya bergerak naik turun. \omterm{def:g10:signals-and-waves:wave}{Gelombang} itu \emph{longitudinal}\index{gelombang
longitudinal} apabila mediumnya bergerak \emph{searah} rambatnya: daerah
lilitan yang termampatkan merambat menyusuri pegas yang didorong, dengan
tiap lilitannya berayun di tempatnya.
\end{definition}

\begin{omfigure}
\begin{tikzpicture}[font=\small]
  \node[left] at (-0.15,2.3) {transversal};
  \draw[omDef, thick, smooth, samples=100, domain=0:8.4]
    plot (\x, {2.3 + 0.42*sin(130*\x)});
  \draw[omThm, Stealth-Stealth] (6.92,1.78) -- (6.92,2.82)
    node[right, yshift=-14pt] {gerak};
  \draw[-Stealth, thick] (3.4,3.15) -- (5.0,3.15)
    node[midway, above] {perambatan};
  \node[left] at (-0.15,0) {longitudinal};
  \foreach \i in {0,...,41}
    \draw[omDef, thick]
      ({0.2*\i + 0.26*sin(32*\i)},0.32) -- ({0.2*\i + 0.26*sin(32*\i)},-0.32);
  \node[below] at (3.4,-0.35) {rapatan};
  \node[below] at (6.65,-0.35) {renggangan};
  \draw[omThm, Stealth-Stealth] (0.75,0.55) -- (1.55,0.55)
    node[midway, above] {gerak};
\end{tikzpicture}

{\small \omterm{def:g12:mechanical-waves:transverse}{Transversal} (tali): mediumnya bergerak tegak lurus rambatnya.
\omterm{def:g12:mechanical-waves:transverse}{Longitudinal} (pegas): rapatannya merambat, lilitannya berayun di
tempatnya.}
\end{omfigure}

\begin{remark}[Bunyi bersifat longitudinal]\label{rem:g12:mechanical-waves:sound}
\omterm{def:g10:signals-and-waves:sound}{Bunyi} adalah \omterm{def:g12:mechanical-waves:transverse}{gelombang longitudinal} yang paling khas: udara yang
termampatkan --- bertekanan lebih tinggi (\cref{ch:g10:pressure}) ---
merambat dari sumbernya ke gendang telinga;
\cref{ch:g12:sound-acoustics} dikhususkan untuknya. Gempa bumi
mengirimkan kedua jenisnya: \omterm{def:g10:signals-and-waves:wave}{gelombang} P memampatkan batuannya, \omterm{def:g10:signals-and-waves:wave}{gelombang}
S mengguncangnya ke samping.
\end{remark}

\section{Laju dan waktu tunda}

\begin{definition}[Laju rambat]\label{def:g12:mechanical-waves:speed}
\emph{Laju rambat}\index{laju rambat} sebuah \omterm{def:g10:signals-and-waves:wave}{gelombang} adalah jarak yang
ditempuh gangguannya dibagi waktu yang diperlukannya: $v = d/t$. Laju
itu adalah laju \emph{bentuk} yang merambat, bukan laju mediumnya.
\end{definition}

\begin{omfigure}
\begin{tikzpicture}[font=\small]
  \draw[-Stealth] (-0.4,0) -- (8.9,0) node[below] {$x$ (\unit{m})};
  \foreach \x in {0,1,...,8}
    \draw (\x,0.05) -- (\x,-0.05) node[below] {$\x$};
  \draw[omDef, thick, smooth, samples=90, domain=-0.2:8.6]
    plot (\x, {1.1*exp(-3.5*(\x-1.0)^2)});
  \draw[omThm, thick, dashed, smooth, samples=90, domain=-0.2:8.6]
    plot (\x, {1.1*exp(-3.5*(\x-3.4)^2)});
  \node[omDef, above] at (1.0,1.15) {$t = 0$};
  \node[omThm, above] at (3.75,1.15) {$t = \qty{0.30}{s}$};
  \draw[omProp, Stealth-Stealth] (1.0,1.75) -- (3.4,1.75)
    node[midway, above] {$d = v\,\Delta t$};
\end{tikzpicture}

{\small Sebuah denyut pada tali yang difoto dua kali: puncaknya
berpindah \qty{2.4}{m} dalam \qty{0.30}{s}, sehingga
$v = \qty{8.0}{m/s}$ --- bentuknya merambat, talinya tinggal.}
\end{omfigure}

\begin{proposition}[Mediumnya yang menetapkan lajunya]\label{prop:g12:mechanical-waves:medium}
Untuk gangguan yang kecil, \omterm{def:g12:mechanical-waves:speed}{laju rambatnya} hanya bergantung pada
\emph{mediumnya}, bukan pada bentuk atau ukuran gangguannya. Pada tali
yang teregang, $v$ bertambah bersama \omterm{def:g11:circuits-and-power:voltage}{tegangannya} dan berkurang bersama
massa tiap \omterm{def:g10:orders-of-magnitude:unit}{meternya}: tegang dan ringan berarti cepat, kendur dan berat
berarti lambat. \admitted
\end{proposition}

\begin{remark}[Di mana rumusnya bersemayam]\label{rem:g12:mechanical-waves:ropeformula}
Rumus laju yang tepat untuk tali diturunkan pada jilid Tahun ke-1; tahun
ini fakta percobaan di atas sudah memadai.
\end{remark}

\begin{example}[Laju bunyi]\label{ex:g12:mechanical-waves:speeds}
\omterm{def:g10:signals-and-waves:sound}{Bunyi} merambat \qty{340}{m/s} di udara, \qty{1500}{m/s} di air, dan
\qty{5900}{m/s} di baja --- makin kaku, makin cepat. Dari sinilah adegan
film koboi yang menempelkan telinga pada relnya: bajalah yang lebih
dahulu mengumumkan keretanya.
\end{example}

\begin{definition}[Waktu tunda]\label{def:g12:mechanical-waves:delay}
\omterm{def:g10:signals-and-waves:wave}{Gelombang} berlaju $v$ mencapai sebuah titik pada jarak $d$ sesudah
\emph{waktu tunda}\index{waktu tunda} $\tau = d/v$: dua titik yang
terpisah sejauh $d$ mengulangi gerak yang sama, dengan yang lebih jauh
tertinggal sebesar $\tau$ --- yaitu pengukuran jarak dengan \omterm{rem:g10:signals-and-waves:honest}{gema} pada
\cref{ch:g10:signals-and-waves}, yang dibaca terbalik.
\end{definition}

\begin{proposition}[Superposisi tanpa saling mengganggu]\label{prop:g12:mechanical-waves:superposition}
Di tempat dua \omterm{def:g10:signals-and-waves:wave}{gelombang} bertindihan, simpangan mediumnya pada setiap
saat adalah \emph{jumlah} simpangan yang akan ditimbulkan
masing-masingnya sendirian; sesudah saling melintas, tiap \omterm{def:g10:signals-and-waves:wave}{gelombangnya}
melanjutkan perjalanan tanpa berubah --- diturunkan pada jilid Tahun
ke-2 dari kelinearan persamaan geraknya. \admitted
\end{proposition}

\begin{example}[Denyut yang saling melintas]\label{ex:g12:mechanical-waves:crossing}
Kirimkanlah sebuah gundukan ke atas dari tiap ujung sebuah tali: di
tempat keduanya bertemu, talinya sejenak naik setinggi jumlah kedua
tingginya; lalu muncullah dua gundukan yang utuh --- seperti dua
percakapan yang saling melintas dan tetap sampai kepada pendengarnya.
\end{example}

\section{Gelombang periodik: keperiodikan ganda}

\begin{definition}[Gelombang periodik]\label{def:g12:mechanical-waves:periodic}
Ketika sumbernya mengulangi geraknya dengan \omterm{def:g10:signals-and-waves:period}{periode} $T$ dan \omterm{def:g10:signals-and-waves:frequency}{frekuensi}
$f = 1/T$ (\cref{ch:g10:signals-and-waves}), sumber itu menyuapkan
\emph{gelombang periodik}\index{gelombang periodik} kepada mediumnya:
tiap titiknya mengulangi geraknya sendiri dengan \omterm{def:g10:signals-and-waves:period}{periode} $T$, dan
masing-masing tertinggal dari tetangganya sebesar \omterm{def:g12:mechanical-waves:delay}{waktu tunda} tempuhnya.
\end{definition}

\begin{definition}[Panjang gelombang]\label{def:g12:mechanical-waves:wavelength}
Potret sekejap sebuah \omterm{def:g12:mechanical-waves:periodic}{gelombang periodik} menampilkan pola yang berulang
dalam ruang; \omterm{def:g10:signals-and-waves:period}{periode} ruangnya --- dari puncak ke puncak --- disebut
\emph{panjang gelombang}\index{panjang gelombang} $\lambda$, dalam
\omterm{def:g10:orders-of-magnitude:unit}{meter}. \omterm{def:g12:mechanical-waves:periodic}{Gelombang periodik} bersifat \emph{periodik
ganda}\index{keperiodikan ganda}: berperiode $T$ dalam waktu pada tiap
titiknya, berperiode $\lambda$ dalam ruang pada tiap saatnya.
\end{definition}

\begin{theorem}[Hubungan gelombang]\label{thm:g12:mechanical-waves:relation}
\omterm{def:g12:mechanical-waves:periodic}{Gelombang periodik} berlaju $v$, berperiode $T$ dan berfrekuensi $f$
memiliki \omterm{def:g12:mechanical-waves:wavelength}{panjang gelombang} $\lambda = v\,T = v/f$.
\end{theorem}

\begin{proof}
Selama satu \omterm{def:g10:signals-and-waves:period}{periode} polanya meluncur maju sejauh $v\,T$; tetapi
sumbernya --- dan karena itu tiap titiknya --- lalu kembali ke keadaan
awalnya, sehingga pola yang meluncur itu berimpit dengan pola aslinya:
jarak pengulangannya adalah $v\,T$.
\end{proof}

\begin{omfigure}
\begin{tikzpicture}[font=\small]
  \draw[-Stealth] (-0.25,0) -- (4.35,0) node[below] {$x$};
  \draw[-Stealth] (0,-1.1) -- (0,1.55) node[left] {$y$};
  \draw[omDef, thick, smooth, samples=110, domain=0:4.1]
    plot (\x, {0.8*sin(160*\x)});
  \draw[omProp, Stealth-Stealth] (0.5625,1.05) -- (2.8125,1.05)
    node[midway, above] {$\lambda$};
  \node at (2.05,-1.55) {potret sekejap $y(x)$, $t$ tetap};
  \begin{scope}[xshift=5.6cm]
    \draw[-Stealth] (-0.25,0) -- (4.35,0) node[below] {$t$};
    \draw[-Stealth] (0,-1.1) -- (0,1.55) node[left] {$y$};
    \draw[omThm, thick, smooth, samples=110, domain=0:4.1]
      plot (\x, {0.8*sin(160*\x)});
    \draw[omProp, Stealth-Stealth] (0.5625,1.05) -- (2.8125,1.05)
      node[midway, above] {$T$};
    \node at (2.05,-1.55) {sinyal $y(t)$, $x$ tetap};
  \end{scope}
\end{tikzpicture}

{\small Sinus yang sama dua kali: potret sekejapnya berulang tiap
$\lambda$ dalam ruang, \omterm{def:g10:signals-and-waves:signal}{sinyal} satu titiknya tiap $T$ dalam waktu;
$\lambda = v\,T$.}
\end{omfigure}

\begin{remark}[Frekuensinya milik sumbernya]\label{rem:g12:mechanical-waves:source}
Ketika berpindah dari satu medium ke medium yang lain --- \omterm{def:g10:signals-and-waves:sound}{bunyi} dari
udara ke air --- sebuah \omterm{def:g10:signals-and-waves:wave}{gelombang} mempertahankan \omterm{def:g10:signals-and-waves:frequency}{frekuensinya}, yang
ditentukan sumbernya, sementara lajunya berubah; menurut $\lambda = v/f$
\omterm{def:g12:mechanical-waves:wavelength}{panjang gelombangnya} pun berubah bersamanya.
\end{remark}

\begin{example}[Nada A, di udara dan di bawah air]\label{ex:g12:mechanical-waves:concertA}
Nada A \qty{440}{Hz} pada \cref{ch:g10:signals-and-waves} memiliki
\omterm{def:g12:mechanical-waves:wavelength}{panjang gelombang} $340/440 = \qty{0.77}{m}$ di udara dan
$1500/440 =
\qty{3.4}{m}$ di air: nada yang sama, yang teregang oleh
medium yang lebih cepat.
\end{example}

\begin{definition}[Sefase dan berlawanan fase]\label{def:g12:mechanical-waves:phase}
Dua titik yang terpisah sejumlah bulat \omterm{def:g12:mechanical-waves:wavelength}{panjang gelombang},
$d = k\lambda$, bergerak serupa pada setiap saat: keduanya
\emph{sefase}\index{sefase}. Titik yang terpisah
$d = (k + \tfrac12)\lambda$ selalu bergerak berlawanan --- puncak
melawan lembah: keduanya \emph{berlawanan fase}\index{berlawanan fase}.
\end{definition}

\begin{method}[Membaca keperiodikan gandanya]\label{met:g12:mechanical-waves:reading}
\begin{enumerate}
  \item Pada rekaman $y(t)$ di sebuah titik tetap, bacalah $T$ dari
        puncak ke puncak; $f = 1/T$.
  \item Pada potret sekejap $y(x)$ pada saat tertentu, bacalah $\lambda$
        dari puncak ke puncak.
  \item Simpulkan $v = \lambda/T = \lambda f$; periksalah terhadap
        pengukuran waktu langsung $v = d/\tau$ apabila tersedia.
  \item Jangan pernah mencampuradukkan kedua grafiknya: sumbu yang satu
        membawa sekon, yang lain membawa \omterm{def:g10:orders-of-magnitude:unit}{meter}.
\end{enumerate}
\end{method}

\section{Muka gelombang dan pelemahan}

\begin{definition}[Muka gelombang]\label{def:g12:mechanical-waves:wavefront}
\emph{Muka gelombang}\index{muka gelombang} adalah garis (atau
permukaan) berisi titik yang dicapai \omterm{def:g10:signals-and-waves:wave}{gelombangnya} pada saat yang sama
--- yaitu garis puncak di kolam. Dari sumber titik, muka \omterm{def:g10:signals-and-waves:wave}{gelombangnya}
berupa lingkaran yang melebar; jauh dari sumber mana pun, muka itu
hampir lurus. \omterm{def:g10:signals-and-waves:wave}{Gelombangnya} merambat tegak lurus muka \omterm{def:g10:signals-and-waves:wave}{gelombangnya},
dengan puncak yang berurutan terpisah $\lambda$.
\end{definition}

\begin{definition}[Pelemahan]\label{def:g12:mechanical-waves:attenuation}
Berkurangnya \omterm{def:g10:signals-and-waves:amplitude}{amplitudo} sebuah \omterm{def:g10:signals-and-waves:wave}{gelombang} selagi merambat disebut
\emph{pelemahan}\index{pelemahan}nya. Ada dua sebab yang berjumlah:
\omterm{def:g10:energy-conservation:energy}{energi} \omterm{def:g12:mechanical-waves:wavefront}{muka gelombang} yang melingkar dibagikan pada puncak yang makin
panjang, dan mediumnya menyerap sedikit selama perjalanannya. \omterm{def:g12:mechanical-waves:wavefront}{Muka
gelombang} yang lurus lolos dari sebab yang pertama --- begitulah alun
dari badai yang jauh dapat menyeberangi samudra.
\end{definition}

\begin{omfigure}
\begin{tikzpicture}[font=\small]
  \fill (0,0) circle (2pt);
  \foreach \r/\c in {0.5/100, 1.0/82, 1.5/64, 2.0/48, 2.5/34}
    \draw[thick, omDef!\c!white] (0,0) circle (\r);
  \draw[omProp, Stealth-Stealth] (315:1.5) -- (315:2.0)
    node[midway, below left, inner sep=1pt] {$\lambda$};
  \draw[-Stealth, omThm, thick] (25:2.55) -- (25:3.25)
    node[right] {$v$};
  \foreach \x/\c in {4.9/100, 5.5/90, 6.1/80, 6.7/70}
    \draw[thick, omDef!\c!white] (\x,-1.5) -- (\x,1.5);
  \draw[-Stealth, omThm, thick] (6.9,0) -- (7.6,0) node[right] {$v$};
\end{tikzpicture}

{\small \omterm{def:g12:mechanical-waves:wavefront}{Muka gelombang} yang melingkar, terpisah satu \omterm{def:g12:mechanical-waves:wavelength}{panjang gelombang},
memudar ketika \omterm{def:g10:energy-conservation:energy}{energinya} tersebar pada puncak yang lebih panjang; muka
yang lurus (kanan) memudar dengan lambat.}
\end{omfigure}

\section{Latihan}

\begin{exercise}[$\star$]\label{exo:g12:mechanical-waves:1}
\omterm{def:g10:signals-and-waves:wave}{Gelombang} stadion menempuh lingkaran tribun sepanjang \qty{320}{m}
dalam \qty{40}{s}. (a) Hitunglah lajunya. (b) Adakah penonton yang
berjalan mengelilingi tribunnya? Lalu apa yang berjalan? (c) Tiap
penonton berdiri \qty{1.5}{s} sesudah tetangganya yang berjarak
\qty{12}{m}: periksalah kesesuaiannya.
\end{exercise}

\begin{exercise}[$\star$]\label{exo:g12:mechanical-waves:2}
\omterm{def:g12:mechanical-waves:transverse}{Transversal} atau \omterm{def:g12:mechanical-waves:transverse}{longitudinal}? (a) sebuah denyut pada tali yang
disentakkan naik turun; (b) sebuah rapatan yang dikirim menyusuri pegas;
(c) \omterm{def:g10:signals-and-waves:sound}{bunyi} di udara; (d) riak kolam (permukaannya naik turun tegak,
lingkarannya merambat mendatar); (e) \omterm{def:g10:signals-and-waves:wave}{gelombang} seismik P (memampatkan
searah rambatnya) dan \omterm{def:g10:signals-and-waves:wave}{gelombang} S (mengguncang ke samping).
\end{exercise}

\begin{exercise}[$\star$]\label{exo:g12:mechanical-waves:3}
Guruh sampai kepadamu \qty{4.0}{s} sesudah kilatnya. (a) Berapa jauh
sambarannya ($v = \qty{340}{m/s}$)? (b) Seorang penyelam mendengar \omterm{def:g12:mechanical-waves:delay}{waktu
tunda} yang sama antara sebuah kilat dan dentuman di bawah air
($v = \qty{1500}{m/s}$): berapa jauh sumber itu? (c) Apa yang harus kamu
ketahui sebelum mengubah \omterm{def:g12:mechanical-waves:delay}{waktu tunda} menjadi jarak?
\end{exercise}

\begin{exercise}[$\star$]\label{exo:g12:mechanical-waves:4}
Hitunglah \omterm{def:g12:mechanical-waves:wavelength}{panjang gelombang} \omterm{def:g10:signals-and-waves:sound}{bunyi} \qty{440}{Hz} di udara dan di air.
Ketika \omterm{def:g10:signals-and-waves:sound}{bunyinya} menyeberang dari udara ke air, di antara $f$, $v$,
$\lambda$, mana yang berubah dan mana yang tetap?
\end{exercise}

\begin{exercise}[$\star$]\label{exo:g12:mechanical-waves:5}
Di sebuah kolam, puncak riak yang bersebelahan berjarak \qty{12}{cm} dan
lingkarannya maju \qty{0.30}{m/s}. Tentukan \omterm{def:g12:mechanical-waves:wavelength}{panjang gelombang}, \omterm{def:g10:signals-and-waves:period}{periode}
dan \omterm{def:g10:signals-and-waves:frequency}{frekuensi} \omterm{def:g10:signals-and-waves:wave}{gelombangnya}.
\end{exercise}

\begin{exercise}[$\star\star$]\label{exo:g12:mechanical-waves:6}
Sebuah denyut pada tali yang panjang difoto dua kali, berselang
\qty{0.30}{s}: puncaknya berada di $x = \qty{1.0}{m}$, lalu
\qty{3.4}{m}. (a) Berapa lajunya? (b) Kapan denyut itu mencapai dinding
di $x = \qty{7.0}{m}$? (c) Sehelai pita disimpulkan di
$x = \qty{5.0}{m}$: bagaimana geraknya ketika denyutnya lewat?
\end{exercise}

\begin{exercise}[$\star\star$]\label{exo:g12:mechanical-waves:7}
Dua denyut ke atas, yang tingginya \qty{3.0}{cm} dan \qty{2.0}{cm},
merambat saling mendekat masing-masing dengan \qty{2.0}{m/s}, dengan
puncaknya terpisah \qty{4.0}{m}. (a) Kapan kedua puncaknya berimpit?
(b) Berapa tinggi terbesar talinya pada saat itu? (c) Bagaimana jika
denyut yang kecil menghadap ke bawah? (d) Bagaimana talinya tepat
sesudahnya?
\end{exercise}

\begin{exercise}[$\star\star$]\label{exo:g12:mechanical-waves:8}
(a) Dua tali jemuran memiliki \omterm{def:g11:circuits-and-power:voltage}{tegangan} yang sama; sebuah denyut dikirim
menyusuri masing-masing. Yang mana yang lebih cepat: yang \omterm{def:g11:lenses-and-eye:lens}{tipis} atau
yang tebal? Mengapa? (b) Seorang pemain gitar menegangkan senarnya: apa
yang terjadi pada laju \omterm{def:g10:signals-and-waves:wave}{gelombangnya}? (c) Apakah menyentakkan dua kali
lebih kuat (denyut yang lebih tinggi) membuat denyutnya merambat lebih
cepat?
\end{exercise}

\begin{exercise}[$\star\star$]\label{exo:g12:mechanical-waves:9}
Seorang pekerja memukul rel yang berjarak \qty{1.8}{km}; kamu menyimak
dengan satu telinga menempel pada bajanya. Hitunglah waktu tibanya
melalui relnya (\qty{5900}{m/s}) dan melalui udaranya (\qty{340}{m/s}),
serta selang antara kedua dentumannya.
\end{exercise}

\begin{exercise}[$\star\star$]\label{exo:g12:mechanical-waves:10}
Sebuah tali membawa \omterm{def:g12:mechanical-waves:periodic}{gelombang periodik} berpanjang \omterm{def:g10:signals-and-waves:wave}{gelombang} \qty{24}{cm}.
\omterm{def:g12:mechanical-waves:phase}{Sefase}, \omterm{def:g12:mechanical-waves:phase}{berlawanan fase}, atau bukan keduanya, untuk jarak antartitik
(a) \qty{48}{cm}; (b) \qty{12}{cm}; (c) \qty{36}{cm}; (d) \qty{30}{cm};
(e) \qty{6.0}{cm}?
\end{exercise}

\begin{exercise}[$\star\star$]\label{exo:g12:mechanical-waves:11}
Sebuah gempa mengirimkan \omterm{def:g10:signals-and-waves:wave}{gelombang} P pada \qty{6.0}{km/s} dan \omterm{def:g10:signals-and-waves:wave}{gelombang}
S pada \qty{3.5}{km/s}, keduanya berperiode \qty{2.0}{s}. Hitunglah
kedua \omterm{def:g12:mechanical-waves:wavelength}{panjang gelombangnya}, bandingkan dengan sebuah bangunan, lalu
jelaskan mengapa seluruh permukiman naik dan turun nyaris bersamaan.
\end{exercise}

\begin{exercise}[$\star\star\star$]\label{exo:g12:mechanical-waves:12}
Rekaman tinggi air pada sebuah pelampung menunjukkan puncak tiap
\qty{0.50}{s}; sebuah foto udara menunjukkan puncak yang terpisah
\qty{0.75}{m}. (a) Tentukan $T$ dan $f$. (b) Tentukan $\lambda$.
(c) Simpulkan lajunya. (d) Pelampung itu menunggangi sebuah puncak pada
$t = 0$: kapan berikutnya? (e) Apakah titik yang berjarak \qty{1.875}{m}
\omterm{def:g12:mechanical-waves:phase}{sefase} dengan pelampungnya, \omterm{def:g12:mechanical-waves:phase}{berlawanan fase}, atau bukan keduanya?
\end{exercise}

\begin{exercise}[$\star\star\star$]\label{exo:g12:mechanical-waves:13}
Sebuah pelampung yang berjangkar menyelesaikan $10$ ayunan penuh dalam
\qty{80}{s} pada alun yang teratur, yang puncaknya terpisah \qty{120}{m}.
(a) Tentukan \omterm{def:g10:signals-and-waves:period}{periode} dan \omterm{def:g10:signals-and-waves:frequency}{frekuensinya}. (b) Tentukan laju alunnya, dalam
\unit{m/s} dan \unit{km/h}. (c) Berapa lama sebuah puncak yang kini
berada di pelampung mencapai pantai yang berjarak \qty{900}{m}?
(d) Apakah pelampungnya ikut hanyut ke pantai bersamanya?
\end{exercise}

\begin{exercise}[$\star\star\star$]\label{exo:g12:mechanical-waves:14}
Dua denyut yang serupa, satu menghadap ke atas dan satu ke bawah,
merambat saling mendekat masing-masing dengan \qty{3.0}{m/s}, dengan
pusatnya terpisah \qty{6.0}{m}. (a) Kapan kedua pusatnya berimpit, dan
seperti apa talinya tampak pada saat itu? (b) Apakah talinya lalu diam?
Ke mana \omterm{def:g10:energy-conservation:energy}{energinya} pergi? (c) Apa yang terjadi berikutnya? (d) Mengapa
``rata sejenak'' tidak berarti ``tidak ada apa-apa di sana''?
\end{exercise}

\begin{exercise}[$\star\star\star$]\label{exo:g12:mechanical-waves:15}
Riak sebuah kerikil menyebar di kolam yang tenang. (a) Hitunglah panjang
puncaknya pada $r = \qty{0.50}{m}$ dan $r = \qty{8.0}{m}$; berapa faktor
pertumbuhannya? (b) \omterm{def:g10:energy-conservation:energy}{Energi} puncaknya kira-kira kekal: apa yang terjadi
pada \omterm{def:g10:energy-conservation:energy}{energi} tiap \omterm{def:g10:orders-of-magnitude:unit}{meternya}, dan pada \omterm{def:g10:signals-and-waves:amplitude}{amplitudonya}? (c) Mengapa alun yang
mukanya lurus jauh lebih sedikit melemah? (d) Dan mengapa \omterm{def:g10:signals-and-waves:wave}{gelombang}
stadion sama sekali tidak melemah?
\end{exercise}

\section{Soal: Membaca sebuah seismogram}

\begin{problem}[{Soal akhir pekan --- membaca sebuah seismogram: dua
gelombang berlomba keluar dari sesar yang patah, tiga lingkaran memaku
pusat gempanya, dan gelombang ketiga yang lambat memberi sebuah garis
pantai peringatannya}]
\label{pb:g12:mechanical-waves:1}
Giliran jaga malam di sebuah observatorium seismologi. Pada pukul
09:14:32 jarum stasiun A melonjak: getaran tajam yang cepat, lalu ayunan
yang lebih lambat dan lebih lebar. Sebuah gempa telah mematahkan batuan
di suatu tempat; dua \omterm{def:g10:signals-and-waves:wave}{gelombang} berlomba keluar menembus kerak buminya:
\omterm{def:g10:signals-and-waves:wave}{gelombang} P, berupa rapatan pada $v_P = \qty{6.0}{km/s}$, dan \omterm{def:g10:signals-and-waves:wave}{gelombang}
S, berupa guncangan menyamping pada $v_S = \qty{3.5}{km/s}$.

\medskip
\noindent\textbf{Bagian I --- Dua \omterm{def:g10:signals-and-waves:wave}{gelombang}, satu jam.}
\begin{enumerate}
  \item Dari pemerian di atas, golongkan \omterm{def:g10:signals-and-waves:wave}{gelombang} P dan S sebagai
        \omterm{def:g12:mechanical-waves:transverse}{transversal} atau \omterm{def:g12:mechanical-waves:transverse}{longitudinal}.
  \item \omterm{def:g10:signals-and-waves:wave}{Gelombang} yang mana yang tiba lebih dahulu di setiap stasiun?
        Hitunglah waktu tempuh keduanya sepanjang $d = \qty{120}{km}$.
  \item Tunjukkan bahwa selang tibanya P--S adalah
        $\Delta t = d\,(1/v_S - 1/v_P)$ lalu hitunglah nilainya untuk
        $d = \qty{120}{km}$.
  \item Simpulkan aturan pengukuran jarak milik ahli seismologi
        $d = \dfrac{v_P\,v_S}{v_P - v_S}\,\Delta t$ lalu tunjukkan bahwa
        koefisiennya sekitar \qty{8.4}{km} tiap sekon selangnya.
  \item Mengapa selangnya tumbuh bersama jaraknya --- dan mengapa hal
        itu merupakan hadiah bagi stasiun yang tidak pernah menyaksikan
        gempanya bermula?
\end{enumerate}

\medskip
\noindent\textbf{Bagian II --- Lingkaran di atas peta.}
\begin{enumerate}[resume]
  \item Stasiun A mengukur $\Delta t_A = \qty{25}{s}$. Berapa jauh pusat
        gempanya dari A?
  \item Jelaskan mengapa satu bacaan ini menempatkan pusat gempanya pada
        sebuah lingkaran, bukan pada sebuah titik.
  \item Stasiun B dan C mengukur $\Delta t_B = \qty{15}{s}$ dan
        $\Delta t_C = \qty{32}{s}$: berapa jaraknya ke pusat gempanya?
  \item Jelaskan bagaimana ketiga lingkarannya bersama-sama memaku pusat
        gempanya.
  \item \omterm{def:g10:signals-and-waves:wave}{Gelombang} P mencapai A pada pukul 09:14:32. Hitunglah waktu
        tempuhnya dari pusat gempanya, dan waktu asal gempanya.
\end{enumerate}

\medskip
\noindent\textbf{Bagian III --- \omterm{def:g10:signals-and-waves:wave}{Gelombang} yang menyeberangi samudra.}
Ketiga lingkarannya bertemu di lepas pantai: sesarnya patah di bawah
dasar laut lalu menggerakkan kolom airnya. \omterm{def:g10:signals-and-waves:wave}{Gelombang} panjang di air yang
dalam merambat dengan laju yang ditetapkan kedalamannya --- sekitar
\qty{200}{m/s} di samudra terbuka, sekitar \qty{10}{m/s} di dekat
pantai (diterima di sini; rumusnya diturunkan pada jilid Tahun ke-1).
Sebuah pelabuhan berjarak $D = \qty{600}{km}$ dari pusat gempanya.
\begin{enumerate}[resume]
  \item Ubahlah \qty{200}{m/s} ke \unit{km/h} lalu sebutkan kendaraan
        yang laju jelajahnya sebanding.
  \item Berapa lama tsunaminya mencapai pelabuhan itu?
  \item Seismometer pelabuhannya sendiri merasakan \omterm{def:g10:signals-and-waves:wave}{gelombang} P lebih
        dahulu. Sesudah berapa lama? Jika alarmnya berbunyi pada saat
        itu, berapa waktu peringatan yang diperoleh garis pantainya
        sebelum tsunaminya tiba?
  \item Di laut lepas, \omterm{def:g10:signals-and-waves:period}{periode} tsunaminya sekitar \qty{15}{min}.
        Hitunglah \omterm{def:g12:mechanical-waves:wavelength}{panjang gelombangnya} di sana, lalu jelaskan mengapa
        sebuah kapal tidak menyadarinya.
  \item Ketika mendekati pantai lajunya turun menjadi \qty{10}{m/s}
        sementara \omterm{def:g10:signals-and-waves:period}{periodenya} tetap. Hitunglah \omterm{def:g12:mechanical-waves:wavelength}{panjang gelombangnya} yang
        baru. Apa yang tentu terjadi pada \omterm{def:g10:signals-and-waves:wave}{gelombang} itu ketika ia
        berdesakan?
\end{enumerate}

\medskip
\noindent\textbf{Bagian IV --- \omterm{def:g10:energy-conservation:energy}{Energinya}, dan laporannya.}
Sebuah gempa yang kuat melepaskan sekitar \qty{1e15}{J}, yang dibawa
pergi \omterm{def:g10:signals-and-waves:wave}{gelombangnya}.
\begin{enumerate}[resume]
  \item Hitunglah panjang \omterm{def:g12:mechanical-waves:wavefront}{muka gelombang} yang melingkar pada
        \qty{10}{km} dari pusat gempanya, lalu pada \qty{210}{km};
        dengan faktor berapa panjang itu tumbuh?
  \item Simpulkan mengapa guncangannya melemah bersama jaraknya bahkan
        di dalam kerak bumi yang tanpa rugi --- lalu sebutkan sebab
        kedua yang ditambahkan kerak bumi yang sebenarnya.
  \item Zat cair tidak sanggup melawan geseran. Stasiun di sisi
        seberang merekam \omterm{def:g10:signals-and-waves:wave}{gelombang} P tetapi tidak merekam \omterm{def:g10:signals-and-waves:wave}{gelombang} S:
        apa yang diungkapkan hal itu tentang inti Bumi?
  \item Kereta cepat pada bab \omterm{def:g10:energy-conservation:energy}{energi} tahun sebelumnya
        (\cref{ch:g11:mechanical-energy}) membawa sekitar \qty{1.5e9}{J}
        \omterm{def:g11:mechanical-energy:kinetic}{energi kinetik}: setara berapa kereta semacam itu \qty{1e15}{J}?
  \item Laporannya, dalam dua kalimat: di mana dan kapan gempanya
        terjadi, dan berapa waktu peringatan yang diperoleh
        pelabuhannya? Berikan ketiga bilangannya.
\end{enumerate}
\end{problem}
